\documentclass[UTF8]{ctexart}
\usepackage{xeCJK}
\usepackage{geometry}
\geometry{left=2cm,right=2cm,top=2.5cm,bottom=2.5cm}
\setlength{\headheight}{12.7pt}

\usepackage{titlesec}
\titleformat{\section}
  {\normalfont\large\bfseries}
  {\thesection}
  {1em}
  {}
\titlespacing*{\section}
  {0pt}
  {3.5ex plus 1ex minus .2ex}
  {2.3ex plus .2ex}

\usepackage{amsmath}
\usepackage{amssymb}
\usepackage{amsthm}
\usepackage{enumitem}
\setlist[enumerate]{itemsep=0pt,topsep=0pt,parsep=0pt,leftmargin=0pt,itemindent=2.5em}
\setlist[itemize]{itemsep=0pt,topsep=0pt,parsep=0pt,leftmargin=0pt,itemindent=2.5em}
\usepackage{fancyhdr}
\pagestyle{fancy}
\fancyhf{}
\usepackage{graphicx}
\usepackage{hyperref}
\usepackage{indentfirst}
\usepackage{lmodern}

\begin{document}
\setlength{\abovedisplayskip}{1pt}
\setlength{\belowdisplayskip}{1pt}

\section{乘积结构}

\hypertarget{sec:定义1.1}{}
\noindent\textbf{定义1.1 (乘积结构)}

给定一阶语言$\mathcal{L}$和一族
$\mathcal{L}$-\href{../01 一阶语言的结构/01 一阶语言的结构#nameddest=sec:定义1.1}
{结构} $(A_i)_{i\in I}$, 再给定
$I$上的\href{../../../01 基础集合论/06 滤子/03 超滤#nameddest=sec:定义1.1}{超滤}
$\mathcal{U}$.
在$\displaystyle P=\prod_{i\in I}A_i$上:
\begin{enumerate}
    \item 对任意的常元$c$, 定义$c^P=(c^{A_i})_{i\in I}$,
    \item 对任意的$n$元函数符号$f$, 定义
    \[
        f^P:P^n\to P,\,
        (p_1,\cdots,p_n)\mapsto(f^{A_i}(p_1(i),\cdots,p_n(i)))_{i\in I},
    \]
    \item 对任意的$n$元关系符号$r$, 定义
    \[
        r^P=\{(p_1,\cdots,p_n)\in P^n\mid\exists u\in\mathcal{U},\,
        \forall i\in u,\,(p_1(i),\cdots,p_n(i))\in r^{A_i}\},
    \]
\end{enumerate}
称$P$为$(A_i)_{i\in I}$在$\mathcal{U}$上的\textbf{乘积结构}.

\vspace{10pt}

\hypertarget{sec:定理1.1}{}
\noindent\textbf{定理1.1}

给定一阶语言$\mathcal{L}$, 一族$\mathcal{L}$-环境$(A_i,\sigma_i)_{i\in I}$
和$I$上的超滤$\mathcal{U}$, 记乘积结构为$P$. 再定义
\[
    S:\mathcal{V}\to P,\,x\mapsto(\sigma_i(x))_{i\in I},
\]
作为$P$上的变元赋值. 那么对任意的项$t$有$t^{(P,S)}=(t^{(A_i,\sigma_i)})_{i\in I}$.

\noindent{\kaishu 证明.}

\begin{enumerate}
    \item 对任意的常元$c$,
    $c^{(P,S)}=c^P=(c^{A_i})_{i\in I}=(c^{(A_i,\sigma_i)})_{i\in I}$,
    \item 对任意的变元$x$,
    $x^{(P,S)}=x^S=(x^{\sigma_i})_{i\in I}=(x^{(A_i,\sigma_i)})_{i\in I}$,
    \item 对任意的$n$元函数符号$f$, 假设项$t_1,\cdots,t_n$满足命题,
    那么对$t=ft_1\cdots t_n$有
    \begin{equation*}
        t^{(P,S)}=f^P(t_1^{(P,S)},\cdots,t_n^{(P,S)})
        =(f^{A_i}(t_1^{(A_i,\sigma_i)},\cdots,t_n^{(A_i,\sigma_i)}))_{i\in I}
        =(t^{(A_i,\sigma_i)})_{i\in I}.
    \tag*{\qedsymbol}\end{equation*}
\end{enumerate}





\newpage

\section{超积结构}

\hypertarget{sec:定义2.1}{}
\noindent\textbf{定义2.1 (超积结构)}

给定一阶语言$\mathcal{L}$, 一族$\mathcal{L}$-结构$(A_i)_{i\in I}$和$I$上的
超滤$\mathcal{U}$, 记乘积结构为$P$. 再定义$P$上的等价关系$\sim_{\mathcal{U}}$:
\[
    \forall p,p'\in P:(p\sim_{\mathcal{U}}p'\iff
    \exists u\in\mathcal{U},\,\forall i\in u:p(i)=p'(i)),\\[3pt]
\]
则$\sim_{\mathcal{U}}$是$P$上的$\mathcal{L}$-同余关系, 称商结构
$^{\displaystyle P}\!/_{\displaystyle\sim_{\mathcal{U}}}$为
$(A_i)_{i\in I}$在$\mathcal{U}$上的\textbf{超积结构},
也记作$^{\displaystyle P}\!/_{\displaystyle\mathcal{U}}$.

\noindent{\kaishu 验证.}

\begin{enumerate}
    \item 对任意的$n$元函数符号$f$和
    $p_1\sim_{\mathcal{U}}p_1',\cdots,p_n\sim_{\mathcal{U}}p_n'\in P$,
    依定义存在$u_1,\cdots,u_n\in\mathcal{U}$使得
    \[
        \forall j\leqslant n,\,\forall i\in u_j:(p_j(i)=p_j'(i)),
    \]
    此时定义$\displaystyle u=\bigcap_{j\leqslant n}u_j\in\mathcal{U}$,
    则对任意的$i\in u$有$\forall j\leqslant n:p_j(i)=p_j'(i)$, 从而
    \vspace{3pt}
    \[
        f^P(p_1,\cdots,p_n)(i)=f^{A_i}(p_1(i),\cdots,p_n(i))
        =f^{A_i}(p_1'(i),\cdots,p_n'(i))=f^P(p_1',\cdots,p_n')(i),
    \]
    依定义即得$f^P(p_1,\cdots,p_n)\sim_{\mathcal{U}}f^P(p_1',\cdots,p_n')$.

    \item 对任意的$n$元关系符号$r$和
    $p_1\sim_{\mathcal{U}}p_1',\cdots,p_n\sim_{\mathcal{U}}p_n'\in P$,
    假设$(p_1,\cdots,p_n)\in r^P$.

    \hspace{2.17em}
    首先, 同理存在$u\in\mathcal{U}$使得
    对任意的$i\in u$有$\forall j\leqslant n:p_j(i)=p_j'(i)$;

    \hspace{2.17em}
    其次根据$r^P$的定义, 又存在$u'\in\mathcal{U}$使得
    对任意的$i\in u'$有$(p_1(i),\cdots,p_n(i))\in r^{A_i}$.

    \hspace{2.17em}
    所以对任意的$i\in u\cap u'$就有$(p_1'(i),\cdots,p_n'(i))\in r^{A_i}$,
    从而$(p_1',\cdots,p_n')\in r^P$. \hfill $\qedsymbol$
\end{enumerate}

\vspace{10pt}

超积结构$^{\displaystyle P}\!/_{\displaystyle\mathcal{U}}$的具体解释如下:
\vspace{3pt}
\begin{enumerate}
    \item 对任意的常元$c$, $c^{P/\mathcal{U}}=[(c^{A_i})_{i\in I}]$,
    \item 对任意的$n$元函数符号$f$和$s_1,\cdots,s_n\in\,\!
    ^{\displaystyle P}\!/_{\displaystyle\mathcal{U}}$,
    \[
        \forall p_1\in s_1,\,\cdots,\,\forall p_n\in s_n:
        f^{P/\mathcal{U}}(s_1,\cdots,s_n)=
        [(f^{A_i}(p_1(i),\cdots,p_n(i)))_{i\in I}],
    \]
    \item 对任意的$n$元关系符号$r$和$s_1,\cdots,s_n\in\,\!
    ^{\displaystyle P}\!/_{\displaystyle\mathcal{U}}$,
    \[
        \forall p_1\in s_1,\,\cdots,\,\forall p_n\in s_n:
        ((s_1,\cdots,s_n)\in r^{P/\mathcal{U}}\leftrightarrow\exists u\in\mathcal{U},\,
        \forall i\in u:(p_1(i),\cdots,p_n(i))\in r^{A_i}).
    \]
\end{enumerate}

\vspace{10pt}

\hypertarget{sec:定理2.1}{}
\noindent\textbf{定理2.1}

给定一阶语言$\mathcal{L}$, 一族$\mathcal{L}$-环境$(A_i,\sigma_i)_{i\in I}$
和$I$上的超滤$\mathcal{U}$,
记超积结构为$^{\displaystyle P}\!/_{\displaystyle\mathcal{U}}$. 再定义
\[
    ^{\displaystyle S}\!/_{\displaystyle\mathcal{U}}:
    \mathcal{V}\to\,^{\displaystyle P}\!/_{\displaystyle\mathcal{U}},\,
    x\mapsto[(\sigma_i(x))_{i\in I}],\\[3pt]
\]
作为$^{\displaystyle P}\!/_{\displaystyle\mathcal{U}}$上的变元赋值.
那么对任意的项$t$有$t^{(P/\mathcal{U},S/\mathcal{U})}
=[(t^{(A_i,\sigma_i)})_{i\in I}]$.

\vspace{10pt}

\hypertarget{sec:定义2.2}{}
\noindent\textbf{定义2.2 (超幂)}

给定一阶语言$\mathcal{L}$, $\mathcal{L}$-结构$A$, 非空集合$I$
和$I$上的超滤$\mathcal{U}$.

称$(A)_{i\in I}$在$\mathcal{U}$上的超积结构为$A$对$\mathcal{U}$的\textbf{超幂},
记作$A^\mathcal{U}$.

\end{document}
